\chapter{Introducere}

\section{Contextul proiectului}

Testarea automată a aplicațiilor web reprezintă o componentă esențială în ciclul de dezvoltare
al oricărui sistem software modern. Pe măsură ce aplicațiile cresc în complexitate, testarea
manuală devine insuficientă atât ca acoperire, cât și ca viteză de execuție.

Acest proiect are ca scop implementarea și compararea a două suite de teste automate E2E
pentru aceeași aplicație țintă, folosind două framework-uri diferite: \textbf{Playwright}
(Python) și \textbf{Cypress} (TypeScript). Prin implementarea duală a acelorași scenarii de
testare, proiectul oferă o perspectivă practică asupra diferențelor dintre cele două
instrumente din punct de vedere al capabilităților tehnice și limitărilor fiecăruia.

\section{Aplicația testată --- NopCommerce}

\textbf{NopCommerce} este o platformă de e-commerce open-source, construită pe ASP.NET Core,
utilizată pe scară largă atât în mediul academic, cât și în producție. Oferă un flux complet
de cumpărături: navigare prin catalog, filtrare produse, gestionarea coșului și listei de
dorințe, înregistrare cont, autentificare, checkout în mai mulți pași și procesare plăți.

Aceastǎ platforma reprezinta un candidat ideal pentru testare E2E, deoarece permite
acoperirea unor scenarii realiste și complexe: interacțiuni asincrone (AJAX), formulare
multi-pas, calcul dinamic de prețuri, gestionarea sesiunii și persistența datelor între
sesiuni.

Testele din acest proiect au fost rulate local, pe o instanță NopCommerce găzduită prin Docker. Instanța publică oficială blochează sesiunile de browser automatizate prin protecția Cloudflare, care detectează traficul generat de framework-uri și returnează erori înainte ca testele să poată interacționa cu aplicația. Framework-ul Playwright a folosit în plus biblioteca \textbf{playwright-stealth} pentru a masca semnăturile de automatizare atunci când instanța publică era accesată.

\section{Obiectivele proiectului}

Proiectul urmărește trei obiective principale:

\begin{enumerate}
    \item \textbf{Definirea și documentarea scenariilor de testare} --- 6 workflow-uri E2E
          complete, reprezentând user flow-urile reale pe o platformă
          de e-commerce.
    \item \textbf{Implementarea scenariilor în Playwright și Cypress} --- același set de
          scenarii este implementat în ambele framework-uri, permițând o comparație directă
          a abordărilor.
    \item \textbf{Analiza comparativă} --- evaluarea celor două instrumente din perspectiva
          configurării, scrierii testelor, depanării și a capabilităților tehnice avansate.
\end{enumerate}

\section{Configurarea mediului de testare}

Aplicația NopCommerce rulează local prin Docker Compose, folosind fișierul
\texttt{docker-compose.nopcommerce.yml}. Acesta orchestrează două servicii:
\textbf{nopcommerce-db} (SQL Server 2022, portul 1433) și \textbf{nopcommerce-web}
(aplicația NopCommerce, portul 8080), cu un \textit{health check} care asigură că
baza de date este funcțională înainte de pornirea aplicației. Pornirea mediului se
face cu:

\begin{lstlisting}[language=bash]
docker compose -f docker-compose.nopcommerce.yml up -d
\end{lstlisting}

Framework-ul Playwright este implementat în \textbf{Python}. Dependențele se instalează
și browserele se descarcă cu:

\begin{lstlisting}[language=bash]
pip install -r requirements.txt
playwright install
\end{lstlisting}

Testele se rulează cu:

\begin{lstlisting}[language=bash]
pytest test_workflows.py
\end{lstlisting}

Framework-ul Cypress este implementat în \textbf{TypeScript} (modul strict). Dependențele
se instalează cu:

\begin{lstlisting}[language=bash]
npm install
\end{lstlisting}

Testele se rulează headless cu \texttt{npm test} sau în modul interactiv cu
\texttt{npm run cy:open}.
